% 06 - REINITIALISATIONS ET DEFAUTS
% ==============================================================
\sectionintro{06 · Robustesse}{Réinitialisations et gestion des défauts}{Les commandes sont neutralisées avant toute reprise ; l'appui long sert également d'acquittement.}

\noindent
\begin{minipage}[t]{0.355\textwidth}\vspace{0pt}
  \subsection{Réinitialisations}
\begin{lstlisting}[style=stcode]
bResetGlobal := FALSE;

FOR i := 1 TO 8 DO
    IF aResetBac[i] THEN
        aPieceBac[i] := 0; // reset bac
    END_IF

    aTempoReset[i](
        IN := aResetBac[i],
        PT := T#10S);

    IF aTempoReset[i].Q THEN
        bResetGlobal := TRUE;
    END_IF
END_FOR
\end{lstlisting}

  \begin{primarycard}
    \textbf{Effet de l'appui long}\par\small
    Le maintien pendant 10 s de l'un des huit boutons remet à zéro les huit compteurs, les deux commandes et \texttt{nImp}. Le défaut est acquitté, le temporisateur du vérin est réinitialisé et le mode repasse à \texttt{INIT}.
  \end{primarycard}
\end{minipage}\hfill
\begin{minipage}[t]{0.605\textwidth}\vspace{0pt}
  \subsection{Conditions de défaut présentes}
  \begin{tabularx}{\linewidth}{@{}L{37mm}Y@{}}
    \toprule
    \textbf{Condition} & \textbf{Réaction}\\
    \midrule
    \texttt{nPieceBac <= 0} & défaut immédiat, temporisateur arrêté et commandes neutralisées\\
    \texttt{tVerin <= T\#0MS} & défaut immédiat, temporisateur arrêté et commandes neutralisées\\
    \texttt{nBac} hors de l'intervalle 1 à 8 & blocage avant accès au tableau et passage en \texttt{DEFAUT}\\
    état cycle incohérent & commandes à \texttt{FALSE} et défaut immédiat\\
    mode incohérent & passage contrôlé en \texttt{DEFAUT}\\
    mode \texttt{DEFAUT} & commandes à \texttt{FALSE}, temporisateur réinitialisé\\
    appui de 10 s sur un bouton & remise à zéro globale, acquittement et retour par \texttt{INIT}\\
    \bottomrule
  \end{tabularx}

  \vspace{3mm}
  \begin{center}
  \begin{tikzpicture}[node distance=9mm,font=\scriptsize]
    \node[draw=ekteacher!35,fill=ekfault,rounded corners=2mm,minimum width=28mm,minimum height=12mm,align=center,text=ekteacher] (d) {DEFAUT\\commandes à \texttt{FALSE}};
    \node[draw=ekline,fill=white,rounded corners=2mm,minimum width=32mm,minimum height=12mm,align=center,right=of d] (c) {cause corrigée\\appui 10 s};
    \node[draw=ekprimary!40,fill=eksoft,rounded corners=2mm,minimum width=28mm,minimum height=12mm,align=center,right=of c] (i) {INIT\\attente zéro};
    \draw[flow] (d)--(c);
    \draw[flow] (c)--(i);
  \end{tikzpicture}
  \end{center}

  \begin{notebox}[colback=ekinfo,colframe=blue!25]{Portée réelle}
    Cette version gère des incohérences logicielles et de paramétrage. Elle ne réalise pas encore de diagnostic mécanique détaillé ni de fonction de sécurité machine.
  \end{notebox}
\end{minipage}

\newpage

% ==============================================================
